\section{Preparing for Battle}

There are several preparations that must be completed before the battle begins. These are described below.

\subsection{Selecting the Scenario}

First, the scenario to be played must be selected. Scenarios are not described in this rulebook but are instead provided in a separate document. In addition to selecting the scenario, the number of points available to each army must also be determined.

\subsection{Building an Army List}

To play NetEpic, you must create an army list. Each Codex contains all the relevant information needed to create an army.

\subsubsection{Points Limit}

Armies are created by purchasing formations up to the agreed points limit. A typical army contains between 3,000 and 5,000 points per player. Larger battles are possible, but games exceeding 10,000 points will take longer than half a day to complete.

\subsubsection{Formations}

An army consists of formations representing its organisation and constituent elements. There are several types of formation:

\begin{description}
    \item[Mandatory Formation:] Some armies are required to include specific formations in their army list.

    \item[Company Formation:] Company formations, often referred to simply as ``companies'', form the core of the army. You must purchase at least one company formation for each army list. Companies consist of one or more detachments and allow you to include other types of formation. \relax

    \item[Special Formation:] Special formations represent units available only in very limited numbers or characters. You may purchase a maximum of one special formation for each company formation in your army list. Some armies may include one additional special formation.

    \item[Support Formation:] Support formations represent the most common troops. Each support formation consists of one detachment. You may purchase up to five support formations for each company formation in your army list.

    \item[Options:] These are upgrades that may be purchased for a formation, such as certain characters, transports, or consumable items. There is no limit to the number of options that may be purchased.

    \item[Limited Formation:] Some formations have additional restrictions. The possible restrictions are:
\end{description}

\begin{itemize}
    \item \textbf{Unique:} The formation may only be included once in the army.
    \item \textbf{One per X points:} The formation is limited to one for every full or partial increment of X points.
    \item \textbf{One per company formation:} The formation is limited to one for each company formation.
\end{itemize}

\subsubsection{Allies}

Some armies may include allies, representing closely aligned forces that frequently fight together.

The available alliances and their restrictions are described in the relevant Codexes. An allied army list must be valid in its own right. It must therefore contain at least one company formation, together with the corresponding support and special formations.

Special rules affect allied units unless the phrase ``from its Codex'' appears in the rule.

\subsection{Setting Up the Battlefield}

To play NetEpic, you need a suitable table on which to set up the battlefield.

\subsubsection{Battlefield Size}

A standard battlefield measures 180~cm by 120~cm. If you are playing with armies that are smaller or larger than average, it may be necessary to reduce or increase the length of the battlefield. Its width should remain 120~cm.

\subsubsection{Terrain Features}

Terrain features should be placed on the battlefield. Ideally, the battlefield should provide dense cover in some areas and clear lines of sight in others.

There is no specific rule governing how many terrain features must be placed, and you may arrange the battlefield however you wish. However, if you are unsure how to arrange the terrain, you can use the following method.

Divide the battlefield into six areas, each measuring 60~cm by 60~cm. Place between two and four terrain areas within each section. A terrain area could represent a forest, an industrial zone, or a ruined district. Terrain should cover approximately 25\% of the battlefield's surface.

\subsection{Pre-Battle Clarifications}

Now that everything is ready for the battle to begin, the players should take a few minutes to agree on the following points:

\begin{description}
    \item[Battlefield:] Review the battlefield together and clarify the characteristics of its terrain features. This includes the class of each terrain feature, the capacity and durability of structures, what each terrain area represents in game terms, and any other relevant details.

    \item[Army Presentation:] Each player now presents their army to their opponent. This includes explaining the army's special rules and abilities, identifying which profile each base uses, particularly when proxies are being used, stating which formations attached characters belong to, identifying which detachments are carried by which transports, and providing any other information the opponent should know.
\end{description}

\subsection{Applying the Scenario}

The chosen scenario's battle preparation instructions must now be followed. This includes choosing battlefield sides, marking the deployment zones, and placing the objectives. Refer to the scenario description for further details.

\subsubsection{Placing Objectives}

In most scenarios, objectives must be placed on the battlefield. An objective cannot be placed within a terrain feature that is impassable to at least one type of base. Objectives therefore cannot be placed in lakes or inside structures.

\subsection{Deploying the Armies}

The armies are now deployed. Some armies may have detachments held in reserve. These detachments are not deployed at this stage but may enter the battlefield later.

\subsubsection{Who Deploys First}

For alternating deployment and infiltration, the player with the greatest number of formations goes first. If both players have the same number of formations, each rolls 1D6. The player with the highest result goes first.

\subsubsection{Placing Fortifications}

Some armies may purchase fortifications. If so, they are deployed at this stage within the army's deployment zone, or according to the \textit{Advanced Fortification} ability if they possess it.

\subsubsection{Deployment}

Hidden deployment is the preferred deployment method. If hidden deployment is not possible, alternating deployment should be used.

If a structure lies partly inside and partly outside a deployment zone, units may only deploy within it if at least half of the structure lies inside the deployment zone.

\paragraph{Hidden Deployment}

This deployment method requires a screen placed across the centre of the battlefield, blocking each player's view of the opposing deployment zone. Both players then deploy their forces simultaneously within their respective deployment zones. Once both players have finished deploying, the screen is removed.

\paragraph{Alternating Deployment}

The players alternate deploying three of their detachments within their respective deployment zones.

\subsubsection{Infiltration}

Formations with the \textit{Infiltration} or \textit{Free Deployment} ability move or deploy after all other detachments have been deployed.